% Self-Evolution Experiment Section

\subsection{SOP Self-Evolution}
\label{sec:exp_self_evolution}

\subsubsection{Objective}

A central design goal of GenericAgent is the ability to \emph{self-evolve}: as the agent repeatedly executes a task, it progressively distills its experience into structured operational procedures (SOPs) and ultimately into executable code, reducing the cognitive load imposed on the LLM with each iteration. This section evaluates whether this three-phase evolution pipeline—natural-language instruction $\rightarrow$ SOP document $\rightarrow$ codified SOP—yields measurable and consistent efficiency gains across both a longitudinal single-task setting and a cross-task benchmark.

\subsubsection{Experimental Setup}

\paragraph{Study I: Longitudinal Single-Task Evolution.}
We selected a representative multi-step research task: performing a deep investigation of the LangChain GitHub repository, which requires (1) filtering the five most recently merged pull requests labeled \texttt{bug}, (2) tracing each PR to its associated issues and extracting the underlying modules involved, and (3) visiting the LangChain documentation site to verify whether each module has a Troubleshooting section. The task was executed for nine consecutive rounds (\#1--\#9). Round \#1 was conducted without any SOP, relying solely on the LLM's in-context reasoning. Rounds \#2--\#5 operated under an iteratively refined natural-language SOP that was updated after each run. Rounds \#6--\#9 used a fully codified SOP in which the execution logic was encapsulated as deterministic code, with the LLM invoked only for dynamic decision points. For each round we recorded wall-clock execution time, number of LLM calls, and total token consumption (input, output, and cache-creation tokens).

\paragraph{Study II: Cross-Task Benchmark.}
To assess the generalizability of the SOP mechanism, we constructed a benchmark of eight real-world web tasks spanning four categories: (A) information retrieval and structured extraction, (B) goal-directed multi-step operations, (C) open-ended environment exploration with constrained decision-making, and (D) long-horizon tasks with error recovery. Each task was executed three times independently by both a baseline agent without SOP memory (OC) and GenericAgent with SOP memory (GA), and the mean token consumption was used as the primary metric.

\subsubsection{Results and Analysis}

\paragraph{Three-Phase Efficiency Progression (Study I).}

Table~\ref{tab:self_evo_three_phase} summarizes the aggregate statistics for each phase. In the initial run (Round \#1), the agent required 32 LLM calls, consumed 222,203 tokens in total, and took 7 minutes 30 seconds to complete the task. During the SOP optimization phase (Rounds \#2--\#5), the average token consumption dropped to 50,865—a 77.1\% reduction relative to Round \#1—while the average number of LLM calls fell to 9. Upon entering the codified SOP phase (Rounds \#6--\#9), performance stabilized at an average of 23,619 tokens and 5.25 LLM calls per run, representing an 89.4\% reduction in token consumption and a 6.4$\times$ reduction in LLM invocations compared to the baseline.

\begin{table}[h]
\centering
\caption{Three-phase summary of Study I (LangChain GitHub research task, 9 rounds).}
\label{tab:self_evo_three_phase}
\begin{tabular}{lcccc}
\toprule
Phase & Rounds & Avg.\ Time & Avg.\ LLM Calls & Avg.\ Total Tokens \\
\midrule
Initial run (no SOP)  & \#1       & 7m 30s & 32   & 222,203 \\
SOP optimization      & \#2--\#5  & 3m 08s & 9    & 50,865  \\
Codified SOP          & \#6--\#9  & 1m 50s & 5.25 & 23,619  \\
\bottomrule
\end{tabular}
\end{table}

The full per-round data are presented in Table~\ref{tab:self_evo_full}. Several patterns merit attention. First, the SOP optimization phase exhibits an \emph{oscillating convergence} pattern in token consumption (66k $\rightarrow$ 50k $\rightarrow$ 52k $\rightarrow$ 36k across Rounds \#2--\#5), reflecting the agent's iterative adaptation to edge cases encountered during execution. Second, once the codified SOP phase begins, token consumption stabilizes within a narrow band of $23\text{k} \pm 1.2\text{k}$ (standard deviation), indicating that the system has reached a \emph{convergent state} in which execution cost is highly predictable and no longer sensitive to task-instance variation. By contrast, a system without SOP memory would restart from scratch on every execution, maintaining a constant cost profile equivalent to Round \#1 indefinitely.

\begin{table}[h]
\centering
\caption{Per-round statistics for Study I. Cache-read tokens dominate Round \#1 due to extensive exploratory context construction.}
\label{tab:self_evo_full}
\resizebox{\textwidth}{!}{%
\begin{tabular}{llccccccc}
\toprule
Round & Phase & Time & LLM Calls & Input Tok. & Output Tok. & Cache-Create & Cache-Read & Total Tok. \\
\midrule
\#1 & Initial run  & 7m 30s & 32 & 15,581 & 7,647  & 15,600 & 183,375 & 222,203 \\
\#2 & SOP optim.   & 4m 19s & 12 & 5,045  & 4,860  & 4,553  & 51,883  & 66,341  \\
\#3 & SOP optim.   & 2m 53s & 8  & 2,538  & 4,598  & 3,155  & 39,534  & 49,825  \\
\#4 & SOP optim.   & 2m 29s & 9  & 3,265  & 3,682  & 3,435  & 41,376  & 51,758  \\
\#5 & SOP optim.   & 2m 50s & 7  & 2,568  & 3,276  & 2,424  & 27,268  & 35,536  \\
\#6 & Codified SOP & 2m 24s & 6  & 1,855  & 1,064  & 1,930  & 20,913  & 25,762  \\
\#7 & Codified SOP & 1m 41s & 5  & 1,931  & 1,126  & 1,708  & 18,249  & 23,014  \\
\#8 & Codified SOP & 1m 35s & 5  & 1,884  & 990    & 1,586  & 18,229  & 22,689  \\
\#9 & Codified SOP & 1m 38s & 5  & 1,323  & 994    & 1,659  & 19,034  & 23,010  \\
\bottomrule
\end{tabular}%
}
\end{table}

Examining the token composition across phases reveals a deeper mechanism. In Round \#1, cache-creation tokens account for 61.4\% of total consumption, reflecting the large volume of novel context the agent must construct during exploratory execution. By Round \#9, this proportion falls to 7.2\%, indicating that the codified SOP has effectively pre-encoded the task's knowledge requirements, leaving almost nothing new for the LLM to ``learn'' at runtime. Correspondingly, the reduction in LLM call count (32 $\rightarrow$ 5, a 84.4\% decrease) is the primary driver of overall efficiency gains: each eliminated call removes an entire \emph{understand--reason--generate} cycle, and the calls that remain in Round \#9 are almost exclusively productive invocations rather than exploratory or error-recovery attempts.

A noteworthy secondary observation is that per-call token efficiency follows a U-shaped trajectory across the three phases: 8,737 tok/call in Round \#1, rising to 9,275 tok/call in the SOP phase (Round \#2 representative), then falling to 5,779 tok/call in the codified phase. The transient increase during the SOP phase reflects an \emph{adaptation overhead}: the agent must load and interpret the SOP document at each call, temporarily increasing per-call context size. Codification eliminates this overhead by embedding all logic into executable structures, so the LLM only handles minimal dynamic inputs.

\paragraph{Self-Purification of the Memory System.}

An important property observed across the nine rounds is that failure experiences from Round \#1—including GitHub API rate-limit errors, incorrect PR filtering logic, and documentation page parsing failures—did not propagate into the SOP or codified implementations. This \emph{self-purification} behavior stems from the Action-Verified-Only principle governing SOP construction: only outcomes of successful tool calls are written into the SOP, while failed attempts, unverified assumptions, and intermediate error states are discarded. As a result, the codified SOP contains exclusively verified, deterministic operations, which explains why Rounds \#6--\#9 achieve near-zero failure rates despite the complexity of the original task.

\paragraph{Cross-Task Generalization (Study II).}

Table~\ref{tab:self_evo_benchmark} reports the token consumption of OC and GA across the eight benchmark tasks. GA achieves an overall token saving rate of 79.3\% (OC mean: 8,592k tokens; GA mean: 1,780k tokens), with per-task saving rates ranging from 61.0\% to 92.4\%. Crucially, GA outperforms OC on every single task, demonstrating that the efficiency gains are not an artifact of task-specific overfitting but reflect a \emph{generalizable} mechanism.

\begin{table}[h]
\centering
\caption{Token consumption comparison between OC (no SOP) and GA (with SOP) across 8 benchmark tasks (Study II). Each condition is run 3 times; mean values are reported in units of $10^3$ tokens.}
\label{tab:self_evo_benchmark}
\resizebox{\textwidth}{!}{%
\begin{tabular}{llccccccccc}
\toprule
Task & Description & OC$_1$ & OC$_2$ & OC$_3$ & OC mean & GA$_1$ & GA$_2$ & GA$_3$ & GA mean & Saving \\
\midrule
A1 & E-commerce price comparison  & 5,330 & 660   & 1,080 & 2,357 & 1,070 & 214 & 302 & 529 & 77.6\% \\
A2 & Public company report extraction & 1,340 & 466   & 584   & 797   & 183   & 64  & 43  & 97  & 87.9\% \\
B1 & News subscription flow        & 1,080 & 814   & 649   & 848   & 489   & 241 & 262 & 331 & 61.0\% \\
B2 & Flight query export           & 1,370 & 2,330 & 2,130 & 1,943 & 489   & 301 & 344 & 378 & 80.6\% \\
C1 & Government service navigation & 775   & 672   & 442   & 630   & 319   & 130 & 144 & 198 & 68.6\% \\
C2 & University course archiving   & 320   & 353   & 259   & 310   & 214   & 55  & 69  & 113 & 63.7\% \\
D1 & Cross-page state transfer     & 1,756 & 1,459 & 1,108 & 1,441 & 174   & 80  & 75  & 110 & 92.4\% \\
D2 & Dynamic financial data crawl  & 275   & 257   & 267   & 266   & 25    & 31  & 23  & 26  & 90.1\% \\
\midrule
\multicolumn{5}{l}{Overall (sum of means)} & 8,592 & & & & 1,780 & \textbf{79.3\%} \\
\bottomrule
\end{tabular}%
}
\end{table}

Category-level analysis reveals a positive correlation between task complexity and SOP benefit. Long-horizon tasks with error recovery (Category D) yield the highest average saving of 92.0\%, followed by information retrieval tasks (Category A, 80.2\%), goal-directed multi-step tasks (Category B, 74.6\%), and open-ended exploration tasks (Category C, 67.0\%). The exceptional gains in Category D are attributable to SOP's \emph{path compression} effect: tasks involving multi-step state transfer and exception handling impose heavy trial-and-error costs on OC, whereas GA's SOP encodes the verified success path directly, bypassing redundant exploration entirely. The comparatively modest gains in Category C reflect the sensitivity of open-ended navigation tasks to environmental change—when page structures deviate from SOP expectations, partial re-exploration is required—pointing to dynamic SOP adaptation as a direction for future improvement.

Within the GA group, a consistent \emph{cold-start} pattern is observed: the first execution of each task consumes substantially more tokens than subsequent runs (e.g., A2: 183k $\rightarrow$ 64k $\rightarrow$ 43k; D1: 174k $\rightarrow$ 80k $\rightarrow$ 75k), as the agent must map general SOP knowledge onto the specific page structure and interaction flow of the target environment. This one-time adaptation cost is amortized rapidly, and subsequent executions converge to a stable low-cost regime. In contrast, OC shows no convergence trend across its three runs (e.g., B2: 1,370k $\rightarrow$ 2,330k $\rightarrow$ 2,130k), confirming that without memory, each execution is effectively independent.

Taken together, Study I and Study II provide complementary evidence for the self-evolution mechanism: Study I demonstrates that efficiency improves \emph{within} a single task over successive iterations (longitudinal depth, up to 89.6\% token reduction), while Study II demonstrates that SOP-encoded knowledge transfers \emph{across} diverse task types (cross-task breadth, 79.3\% average saving). This depth-by-breadth validation establishes SOP self-evolution as a robust and generalizable capability of GenericAgent.
